\chapter{Menyiapkan Lingkungan dan Membaca Sumber Data}
\label{modul:lingkungan-sumber-data}

\section{Identitas Modul}

\begin{identitasmodul}
\textbf{Sumber materi}    & Pertemuan 2 \\
\textbf{CPMK}             & CPMK-072 \\
\textbf{Minggu pelaksanaan} & Minggu ke-3 \\
\textbf{Durasi tatap muka} & 120 menit \\
\textbf{Estimasi tugas}   & 3--4 jam di luar sesi \\
\textbf{Moda kerja}       & Individual \\
\textbf{Perangkat}        & Docker Compose, PostgreSQL 17, DBeaver, dan Git \\
\textbf{Dataset}          & Basis data operasional NusaMart, ukuran kecil \\
\textbf{Skrip pendukung}  & \berkas{sql/modul-01/01-inventarisasi.sql} dan \berkas{sql/modul-01/check.sql} \\
\textbf{Luaran}           & \berkas{inventaris-sumber.md}, \berkas{desain-konseptual-v1.md}, dan \berkas{bus-matrix.md} \\
\end{identitasmodul}

\slideref{2}

\section{Capaian dan Indikator Keberhasilan}

Setelah praktikum, mahasiswa diharapkan mampu menjalankan lingkungan bersama,
menginventarisasi sumber data berdasarkan bukti, serta merumuskan proses
bisnis, grain, fakta, dan dimensi awal. Modul dinyatakan berhasil apabila kedua
basis data dapat diakses dan grain dinyatakan dalam satu kalimat bisnis.

\begin{capaian}
Pada akhir modul ini mahasiswa mampu:
\begin{enumerate}[leftmargin=1.4em]
  \item menjalankan lingkungan praktikum dari \berkas{docker-compose.yml} dan
        membuktikan bahwa basis data sumber maupun gudang dapat diakses;
  \item menyusun inventarisasi sumber yang memuat jumlah baris hasil query,
        kunci, relasi antartabel, dan kolom bermasalah;
  \item membedakan tujuan sistem operasional dari tujuan gudang data, serta
        menjelaskan mengapa keduanya tidak dirancang dengan cara yang sama;
  \item merumuskan satu proses bisnis beserta grain, kandidat measure, dan
        kandidat dimensi dalam bentuk dokumen desain konseptual.
\end{enumerate}
\end{capaian}

Modul ini sengaja tidak meminta mahasiswa merancang tabel apa pun. Yang diminta
adalah membaca sumber sampai tuntas. Kesalahan desain yang paling mahal pada
gudang data hampir selalu berakar pada sumber yang tidak dipahami: grain yang
salah ditebak, kolom yang dikira unik ternyata tidak, atau baris yang dikira
transaksi sah ternyata sudah dibatalkan.

\section{Prasyarat dan Persiapan}

\subsection{Prasyarat}

Mahasiswa perlu menguasai hal berikut sebelum masuk laboratorium:

\begin{itemize}
  \item SQL dasar sampai \perintah{JOIN} dan \perintah{GROUP BY}, termasuk
        fungsi agregat \perintah{COUNT}, \perintah{SUM}, dan
        \perintah{COUNT(DISTINCT \ldots)};
  \item konsep basis data relasional: primary key, foreign key, kardinalitas
        relasi, dan normalisasi sampai bentuk normal ketiga;
  \item perintah dasar terminal --- berpindah direktori, menjalankan program,
        dan membaca keluaran galat;
  \item Git tingkat dasar: \perintah{clone}, \perintah{add}, \perintah{commit},
        dan \perintah{push}.
\end{itemize}

\subsection{Pre-lab}

Pekerjaan berikut diselesaikan \emph{sebelum} sesi laboratorium. Tanpa ini,
120 menit akan habis untuk pemasangan perangkat lunak, bukan untuk menganalisis
data.

\begin{enumerate}[leftmargin=1.4em]
  \item Pasang Docker Desktop, DBeaver Community, dan Git. Pastikan
        \perintah{docker --version} dan \perintah{git --version} memberi
        keluaran versi, bukan pesan \emph{command not found}.
  \item Unduh berkas starter dari repositori kelas: \berkas{docker-compose.yml},
        direktori \berkas{seed/}, dan \berkas{README.md}.
  \item Jalankan \perintah{docker compose pull} di rumah. Unduhan image
        PostgreSQL berukuran ratusan megabyte dan tidak boleh membebani jaringan
        laboratorium secara serentak.
  \item Buat repositori pribadi \berkas{praktikum-pgd-<nim>} dan buat direktori
        \berkas{modul-01-setup/} di dalamnya.
\end{enumerate}

\begin{peringatan}
Port 5432 sering sudah dipakai oleh PostgreSQL yang terpasang langsung di
laptop. Bila demikian, \perintah{docker compose up} akan gagal dengan pesan
\emph{port is already allocated}. Penyelesaiannya ada pada
Seksi~\ref{sec:m1-troubleshooting}; periksa hal ini saat pre-lab, bukan saat
laboratorium berjalan.
\end{peringatan}

\section{Konsep Dasar}

\subsection{Sumber Operasional dan Gudang Data}

Sistem operasional NusaMart dibangun untuk satu tujuan: mencatat transaksi
dengan benar, cepat, dan tanpa duplikasi. Gudang data dibangun untuk tujuan yang
berbeda: menjawab pertanyaan tentang banyak transaksi sekaligus. Perbedaan
tujuan ini melahirkan perbedaan rancangan yang sistematis, sebagaimana dirangkum
pada Tabel~\ref{tab:m1-oltp-dw}.

\begin{table}[htbp]
\centering
\small
\caption{Perbedaan rancangan sistem operasional dan gudang data}
\label{tab:m1-oltp-dw}
\begin{tabular}{@{}p{0.24\textwidth}p{0.34\textwidth}p{0.34\textwidth}@{}}
\toprule
\textbf{Aspek} & \textbf{Sistem operasional (OLTP)} & \textbf{Gudang data} \\
\midrule
Pertanyaan khas & Berapa total belanja struk nomor 4471? &
  Bagaimana tren penjualan kategori minuman per provinsi selama dua tahun? \\
Satuan kerja & Satu baris pada satu waktu & Jutaan baris sekaligus \\
Rancangan tabel & Ternormalisasi, meminimalkan redundansi &
  Terdenormalisasi, meminimalkan jumlah join \\
Perlakuan histori & Nilai lama ditimpa & Nilai lama dipertahankan \\
Ukuran keberhasilan & Waktu tanggap per transaksi & Waktu tanggap per query
  analitik \\
\bottomrule
\end{tabular}
\end{table}

Satu konsekuensi perlu ditegaskan sejak awal. Ketika bagian pemasaran mengubah
kategori sebuah produk, sistem operasional cukup menimpa nilai lama karena yang
dibutuhkan hanyalah kategori \emph{saat ini}. Gudang data tidak boleh melakukan
hal yang sama: laporan penjualan tahun lalu harus tetap memakai kategori yang
berlaku tahun lalu. Persoalan inilah yang ditangani pada Modul~3 dengan
\istilah{slowly changing dimension}. Pada modul ini mahasiswa cukup
mengenalinya sebagai risiko dan mencatat kolom mana saja yang berpotensi
berubah \citep{kimball2013}.

\subsection{Proses Bisnis, Grain, Fakta, dan Dimensi}

Perancangan dimensional dimulai dari empat keputusan yang diambil berurutan
\citep{kimball2013}. Urutannya penting: menukar urutan hampir selalu
menghasilkan skema yang tidak dapat menjawab pertanyaan yang diinginkan.

\begin{konsep}{Empat langkah perancangan dimensional}
\begin{enumerate}[leftmargin=1.4em]
  \item \textbf{Pilih proses bisnis.} Satu kegiatan operasional yang
        menghasilkan jejak terukur --- misalnya penjualan di kasir, bukan
        ``penjualan'' sebagai konsep umum.
  \item \textbf{Tetapkan grain.} Satu kalimat yang menjawab: \emph{satu baris
        fakta mewakili apa?} Grain ditetapkan sebelum measure dan dimensi
        dipilih, bukan sesudahnya.
  \item \textbf{Pilih dimensi.} Semua konteks yang menjawab siapa, apa, kapan,
        di mana, dan bagaimana, sepanjang konsisten dengan grain.
  \item \textbf{Pilih measure.} Besaran numerik yang bermakna pada grain
        tersebut dan, sedapat mungkin, bersifat aditif.
\end{enumerate}
\end{konsep}

Untuk NusaMart, kalimat grain yang baik berbunyi:

\begin{quote}
\emph{Satu baris fakta mewakili satu produk pada satu baris struk penjualan di
satu toko pada satu waktu.}
\end{quote}

Perhatikan bahwa kalimat ini tidak menyebut tabel, kolom, atau tipe data. Grain
dinyatakan dalam bahasa bisnis justru agar dapat diperiksa oleh orang yang
paham operasi toko tetapi tidak paham SQL. Bandingkan dengan rumusan yang
lemah, ``satu baris mewakili penjualan'': kalimat itu tidak menjawab apakah
satu baris berarti satu struk, satu produk pada satu struk, atau rekap harian
satu toko. Ketiganya menghasilkan tabel fakta yang sama sekali berbeda.

\begin{catatan}
Grain yang lebih halus selalu dapat diringkas menjadi grain yang lebih kasar,
tetapi tidak sebaliknya. Karena itu, bila ragu, pilih grain paling halus yang
tersedia pada sumber. Rekap harian dapat dihitung dari baris struk; baris struk
tidak dapat dipulihkan dari rekap harian.
\end{catatan}

Dari grain di atas, kandidatnya menjadi jelas: measure berupa kuantitas,
harga satuan, diskon, dan subtotal; dimensi berupa produk, toko, waktu,
pelanggan, promosi, dan kasir.

\subsection{Inventarisasi dan Profiling Sumber}

Inventarisasi sumber adalah dokumen yang menjawab pertanyaan berikut untuk
setiap tabel: berapa barisnya, apa kuncinya, ke mana relasinya, dan kolom mana
yang tidak dapat dipercaya. Dokumen ini disusun dari hasil query, bukan dari
membaca nama kolom.

\begin{konsep}{Enam pemeriksaan minimum per tabel}
\begin{enumerate}[leftmargin=1.4em]
  \item \textbf{Volume.} \perintah{COUNT(*)} sebenarnya --- bukan angka
        perkiraan dari panel statistik DBeaver.
  \item \textbf{Keunikan kunci.} Bandingkan \perintah{COUNT(*)} dengan
        \perintah{COUNT(DISTINCT kunci)}. Bila berbeda, kunci yang diduga
        ternyata bukan kunci.
  \item \textbf{Kelengkapan.} Proporsi \perintah{NULL} per kolom, terutama pada
        kolom yang akan menjadi foreign key dimensi.
  \item \textbf{Rentang nilai.} Nilai minimum dan maksimum kolom numerik dan
        tanggal. Nilai ekstrem biasanya menandai anomali, bukan transaksi luar
        biasa.
  \item \textbf{Kardinalitas.} Jumlah nilai berbeda pada kolom kategorikal.
        Kolom berkardinalitas rendah adalah kandidat atribut dimensi.
  \item \textbf{Integritas relasi.} Jumlah baris anak yang tidak menemukan
        pasangan di tabel induk --- yaitu \istilah{orphan}.
\end{enumerate}
\end{konsep}

Temuan profiling tidak dibuang. Setiap kolom bermasalah yang ditemukan hari ini
akan muncul kembali sebagai keputusan desain pada modul berikutnya: baris
\perintah{unknown} pada Modul~5, standardisasi pada Modul~7, dan uji kualitas
pada Modul~10 \citep{batini2016}.

\section{Studi Kasus dan Protokol Praktikum}

NusaMart menjadi kasus utama sepanjang sepuluh modul. Pada modul ini,
mahasiswa memperlakukan sistem operasional sebagai sumber yang harus dipahami
sebelum skema dimensional dirancang.

NusaMart adalah peritel dengan 240 toko yang tersebar di beberapa provinsi dan
sekitar 180.000 produk aktif. Sistem operasionalnya berjalan pada basis data
\perintah{nusamart\_oltp}, dan seluruh tabelnya berada pada schema
\perintah{nusamart}. Tabel~\ref{tab:m1-skema-sumber} merangkum isinya.

\begin{table}[htbp]
\centering
\small
\caption{Tabel pada basis data operasional NusaMart}
\label{tab:m1-skema-sumber}
\begin{tabular}{@{}p{0.26\textwidth}p{0.66\textwidth}@{}}
\toprule
\textbf{Tabel} & \textbf{Isi} \\
\midrule
\perintah{toko} & Identitas, alamat, kota, provinsi, tipe, dan tanggal buka
  setiap gerai. \\
\perintah{kategori} & Hierarki kategori produk, mengacu pada dirinya sendiri
  melalui \perintah{induk\_kategori\_id}. \\
\perintah{produk} & SKU, nama, merek, satuan, berat, harga jual, harga pokok,
  dan penanda aktif. \\
\perintah{pelanggan} & Data anggota program loyalitas: kode member, nama, kota,
  tanggal daftar, dan segmen. \\
\perintah{pegawai} & Kasir dan penanggung jawab toko. \\
\perintah{transaksi} & Kepala struk: nomor struk, toko, pelanggan, kasir, waktu,
  metode bayar, dan status. \\
\perintah{transaksi\_item} & Baris struk: produk, kuantitas, harga satuan,
  diskon, dan subtotal. \\
\perintah{promosi} & Nama, tipe, serta tanggal mulai dan selesai setiap
  program promosi. \\
\perintah{promosi\_produk} & Produk mana saja yang tercakup oleh satu promosi. \\
\perintah{persediaan\_harian} & Saldo awal, masuk, keluar, dan saldo akhir per
  produk per toko per hari. \\
\perintah{retur} dan \perintah{retur\_item} & Pengembalian barang beserta
  alasannya. \\
\bottomrule
\end{tabular}
\end{table}

\begin{catatanasisten}
Basis data ini dibangkitkan dengan skrip \berkas{seed/generate\_nusamart.py}
yang memakai pustaka \perintah{Faker}. Skrip menanam sejumlah anomali secara
sengaja dan terkendali --- lihat Bagian~C. Jalankan skrip dengan
\perintah{--ukuran kecil} untuk Modul 1--5. Nilai \perintah{--seed} harus sama
untuk seluruh kelas agar angka hasil profiling dapat dibandingkan antarkelompok
saat checkpoint.
\end{catatanasisten}

Selain \perintah{nusamart\_oltp}, lingkungan menyediakan basis data kosong
\perintah{nusamart\_dw} yang akan diisi mulai Modul~2. Dua basis data yang
terpisah sejak awal membuat batas antara sumber dan gudang terasa nyata: tidak
ada query yang dapat menyeberang begitu saja dari satu ke yang lain.

\section{Alur Praktikum 120 Menit}

\begin{alurwaktu}
0--15 & Briefing dan demonstrasi lingkungan oleh asisten. \\
15--95 & Kerja terbimbing: koneksi, pemulihan data, profiling, dan desain awal. \\
95--110 & Checkpoint dengan bukti koneksi dan inventaris sumber. \\
110--120 & Penjelasan tugas luar laboratorium. \\
\end{alurwaktu}

\section{Prosedur Praktikum Terbimbing}

\subsection{Bagian A: Menjalankan Lingkungan}

\langkah{A.1 Menyalakan layanan}{10 menit}

Dari direktori yang memuat \berkas{docker-compose.yml}, jalankan:

\begin{lstlisting}[style=shellgaya]
docker compose up -d
docker compose ps
\end{lstlisting}

Keluaran \perintah{docker compose ps} harus memperlihatkan dua layanan berstatus
\emph{running} atau \emph{healthy}: \perintah{postgres\_source} pada port 5433
dan \perintah{postgres\_dw} pada port 5434. Port sengaja digeser dari 5432 agar
tidak berbenturan dengan PostgreSQL yang mungkin sudah terpasang di laptop.

\langkah{A.2 Menguji koneksi dari baris perintah}{5 menit}

\begin{lstlisting}[style=shellgaya]
docker compose exec postgres_source \
  psql -U praktikum -d nusamart_oltp -c "SELECT version();"
\end{lstlisting}

Simpan keluaran perintah ini; potongan versi PostgreSQL menjadi bukti pertama
pada laporan.

\begin{luaran}
\berkas{modul-01-setup/bukti-lingkungan.txt} --- berisi keluaran
\perintah{docker compose ps} dan \perintah{SELECT version();}.
\end{luaran}

\subsection{Bagian B: Menghubungkan dan Memulihkan Sumber}

\langkah{B.1 Membuat dua koneksi DBeaver}{10 menit}

Buat dua koneksi terpisah dengan parameter pada
Tabel~\ref{tab:m1-koneksi}. Beri nama yang jelas, misalnya
\emph{NusaMart--Sumber} dan \emph{NusaMart--Gudang}. Penamaan yang jelas
mencegah kesalahan yang paling sering terjadi sepanjang semester: menjalankan
DDL gudang pada basis data sumber.

\begin{table}[htbp]
\centering
\small
\caption{Parameter koneksi lingkungan praktikum}
\label{tab:m1-koneksi}
\begin{tabular}{@{}llll@{}}
\toprule
\textbf{Peran} & \textbf{Host} & \textbf{Port} & \textbf{Basis data} \\
\midrule
Sumber operasional & \perintah{localhost} & 5433 & \perintah{nusamart\_oltp} \\
Gudang data        & \perintah{localhost} & 5434 & \perintah{nusamart\_dw} \\
\bottomrule
\end{tabular}

\vspace{.5em}
{\footnotesize Pengguna \perintah{praktikum}; kata sandi tercantum pada
\berkas{.env.example}.}
\end{table}

\langkah{B.2 Memeriksa kelengkapan pemulihan}{5 menit}

\begin{lstlisting}[style=sqlgaya]
SELECT table_name,
       (xpath('/row/c/text()',
              query_to_xml(format('SELECT COUNT(*) AS c FROM nusamart.%I',
                                  table_name),
                           false, true, '')))[1]::text::bigint AS jumlah_baris
FROM   information_schema.tables
WHERE  table_schema = 'nusamart'
ORDER  BY jumlah_baris DESC;
\end{lstlisting}

Query ini menghitung baris seluruh tabel sekaligus. Bandingkan hasilnya dengan
angka acuan pada \berkas{README.md} berkas starter. Bila ada tabel yang kosong
atau hilang, pemulihan belum selesai --- ulangi sebelum melanjutkan.

\langkah{B.3 Membuat diagram E/R sumber}{5 menit}

Pada DBeaver, buka schema \perintah{nusamart}, pilih tab \emph{ER Diagram},
lalu ekspor sebagai PNG ke \berkas{modul-01-setup/erd-sumber.png}. Rapikan tata
letak secukupnya sehingga relasi \perintah{transaksi} ke
\perintah{transaksi\_item} dan \perintah{produk} ke \perintah{kategori} terlihat
jelas.

\subsection{Bagian C: Menginventarisasi Sumber Data}

Bagian ini adalah inti modul dan memperoleh porsi waktu terbesar. Kerjakan
enam pemeriksaan yang telah dibahas pada Seksi Konsep Dasar. Skrip pendukung
tersedia pada \berkas{sql/modul-01/01-inventarisasi.sql}.

\langkah{C.1 Volume dan struktur}{10 menit}

Isi Tabel~\ref{tab:m1-inventaris} untuk seluruh tabel sumber. Kolom
\emph{jumlah baris} wajib berasal dari \perintah{COUNT(*)}.

\begin{table}[htbp]
\centering
\small
\caption{Kerangka inventarisasi sumber yang harus dilengkapi}
\label{tab:m1-inventaris}
\begin{tabular}{@{}p{0.19\textwidth}p{0.16\textwidth}p{0.19\textwidth}p{0.34\textwidth}@{}}
\toprule
\textbf{Tabel} & \textbf{Jumlah baris} & \textbf{Kunci} & \textbf{Catatan} \\
\midrule
\perintah{transaksi} & \ldots & \ldots & \ldots \\
\perintah{transaksi\_item} & \ldots & \ldots & \ldots \\
\ldots & \ldots & \ldots & \ldots \\
\bottomrule
\end{tabular}
\end{table}

\langkah{C.2 Menguji dugaan kunci}{10 menit}

\begin{lstlisting}[style=sqlgaya]
-- Apakah nomor_struk unik secara global?
SELECT COUNT(*)                    AS baris,
       COUNT(DISTINCT nomor_struk) AS struk_unik
FROM   nusamart.transaksi;

-- Jika tidak, apakah unik per toko?
SELECT COUNT(*) AS kombinasi_ganda
FROM (
  SELECT toko_id, nomor_struk
  FROM   nusamart.transaksi
  GROUP  BY toko_id, nomor_struk
  HAVING COUNT(*) > 1
) t;
\end{lstlisting}

Catat temuannya dengan kalimat lengkap, bukan sekadar angka. Contoh rumusan
yang benar: ``\perintah{nomor\_struk} tidak unik secara global; keunikan hanya
berlaku pada pasangan \perintah{(toko\_id, nomor\_struk)}.'' Temuan ini
menentukan apa yang akan menjadi \istilah{degenerate dimension} pada Modul~2.

\langkah{C.3 Kelengkapan dan rentang nilai}{15 menit}

\begin{lstlisting}[style=sqlgaya]
SELECT COUNT(*) AS total,
       COUNT(*) FILTER (WHERE pelanggan_id IS NULL) AS tanpa_pelanggan,
       ROUND(100.0 * COUNT(*) FILTER (WHERE pelanggan_id IS NULL)
             / COUNT(*), 2) AS persen_null,
       MIN(waktu_transaksi) AS paling_awal,
       MAX(waktu_transaksi) AS paling_akhir
FROM   nusamart.transaksi;

SELECT status, COUNT(*) AS jumlah
FROM   nusamart.transaksi
GROUP  BY status
ORDER  BY jumlah DESC;
\end{lstlisting}

\begin{peringatan}
Sebagian besar mahasiswa akan menemukan bahwa kolom \perintah{pelanggan\_id}
bernilai \perintah{NULL} pada proporsi transaksi yang besar, lalu menyimpulkan
bahwa datanya rusak. Datanya tidak rusak. \perintah{NULL} di sini berarti
pembeli tidak menunjukkan kartu anggota --- sebuah fakta bisnis yang sah.
Tuliskan sebagai temuan, bukan sebagai kerusakan. Cara gudang data menanganinya
dibahas pada Modul~5 melalui baris \perintah{unknown}.
\end{peringatan}

\langkah{C.4 Integritas relasi dan anomali}{15 menit}

\begin{lstlisting}[style=sqlgaya]
-- Baris struk yang menunjuk produk tidak dikenal.
SELECT COUNT(*) AS item_yatim
FROM      nusamart.transaksi_item ti
LEFT JOIN nusamart.produk p ON p.produk_id = ti.produk_id
WHERE     p.produk_id IS NULL;

-- Konsistensi aritmetika baris struk.
SELECT COUNT(*) AS subtotal_tidak_cocok
FROM   nusamart.transaksi_item
WHERE  ROUND(kuantitas * harga_satuan - diskon, 2) <> ROUND(subtotal, 2);

-- Nilai yang tidak masuk akal.
SELECT COUNT(*) FILTER (WHERE kuantitas <= 0)     AS kuantitas_tidak_positif,
       COUNT(*) FILTER (WHERE diskon < 0)         AS diskon_negatif,
       COUNT(*) FILTER (WHERE harga_satuan <= 0)  AS harga_tidak_positif
FROM   nusamart.transaksi_item;
\end{lstlisting}

Temukan sekurang-kurangnya \textbf{lima} kolom atau kondisi bermasalah, lalu
untuk masing-masing tuliskan tiga hal: apa temuannya, berapa banyak barisnya,
dan apa dugaan penyebabnya.

\begin{luaran}
\berkas{modul-01-setup/inventaris-sumber.md} --- memuat tabel inventarisasi
lengkap, daftar temuan kualitas data beserta jumlah barisnya, dan seluruh query
yang dipakai untuk memperolehnya.
\end{luaran}

\subsection{Bagian D: Menyusun Desain Konseptual Awal}

\langkah{D.1 Merumuskan proses bisnis dan grain}{10 menit}

Pilih \emph{satu} proses bisnis, lalu tulis pernyataan grainnya dalam satu
kalimat bisnis. Uji kalimat itu dengan tiga pertanyaan:

\begin{enumerate}[leftmargin=1.4em]
  \item Dapatkah seorang manajer toko memahaminya tanpa penjelasan tambahan?
  \item Apakah kalimatnya bebas dari nama tabel dan nama kolom?
  \item Dapatkah Anda menunjuk baris tertentu pada \perintah{transaksi\_item}
        dan berkata, ``ini satu baris fakta''?
\end{enumerate}

\langkah{D.2 Memilih measure dan dimensi}{10 menit}

Susun daftar kandidat measure beserta sifat keteraditifannya, dan daftar
kandidat dimensi beserta atribut yang menarik untuk analisis. Untuk setiap
measure, nyatakan apakah ia aditif, semi-aditif, atau non-aditif terhadap
dimensi waktu --- dan jelaskan alasannya.

\begin{catatan}
Harga satuan adalah contoh yang layak diperdebatkan. Menjumlahkan harga satuan
dari sepuluh baris struk menghasilkan angka yang tidak bermakna, sehingga ia
non-aditif. Namun ia tetap perlu disimpan karena diperlukan untuk menghitung
harga rata-rata tertimbang. Menyimpan measure non-aditif bukan kesalahan;
\emph{menjumlahkannya} yang keliru.
\end{catatan}

\langkah{D.3 Merumuskan pertanyaan bisnis}{5 menit}

Tulis \textbf{empat} pertanyaan bisnis yang seharusnya dapat dijawab oleh
gudang data ini. Pertanyaan harus spesifik dan terukur, misalnya ``kategori
produk mana yang kontribusi penjualannya turun lebih dari 10 persen di
Lampung pada kuartal terakhir?'' --- bukan ``bagaimana kinerja penjualan?''.

\begin{catatan}
Simpan keempat pertanyaan ini baik-baik. Pada Modul~9, dashboard yang Anda
bangun harus menjawabnya. Praktikum ini sengaja dirancang menutup lingkaran:
pertanyaan dirumuskan di modul pertama dan dijawab di modul kesembilan.
\end{catatan}

\begin{luaran}
\berkas{modul-01-setup/desain-konseptual-v1.md} --- memuat proses bisnis
terpilih, pernyataan grain, daftar measure beserta sifat keteraditifannya,
daftar dimensi beserta atributnya, dan empat pertanyaan bisnis.
\end{luaran}

\begin{checkpoint}
Asisten menjalankan \berkas{sql/modul-01/check.sql} pada basis data mahasiswa
dan memeriksa butir berikut:
\begin{ceklis}
  \item kedua basis data terhubung dan \perintah{SELECT version();} memberi
        keluaran;
  \item seluruh tabel \perintah{nusamart} terpulihkan dengan jumlah baris sesuai
        angka acuan;
  \item inventarisasi memuat jumlah baris hasil \perintah{COUNT(*)}, bukan
        perkiraan;
  \item sekurang-kurangnya lima temuan kualitas data tercatat beserta jumlah
        barisnya;
  \item grain dinyatakan dalam satu kalimat bisnis tanpa nama tabel dan kolom;
  \item empat pertanyaan bisnis tertulis dan bersifat terukur.
\end{ceklis}
\end{checkpoint}

\section{Latihan Mandiri di Kelas}

\latihan{Menyelidiki satu relasi yang meragukan}

Tabel \perintah{kategori} mengacu pada dirinya sendiri melalui kolom
\perintah{induk\_kategori\_id}. Selidiki hierarki tersebut dan jawab:

\begin{enumerate}[leftmargin=1.4em]
  \item Berapa tingkat kedalaman maksimum hierarki kategori?
  \item Adakah kategori yang menunjuk induk yang tidak ada?
  \item Adakah lingkaran acuan --- kategori yang menjadi leluhur dirinya
        sendiri?
\end{enumerate}

Gunakan \perintah{WITH RECURSIVE} untuk menelusuri hierarki. Tuliskan
konsekuensinya bagi rancangan \perintah{dim\_produk}: bila kedalaman hierarki
tidak seragam, bagaimana kategori akan diratakan menjadi kolom-kolom dimensi?

\latihan{Menakar satu anomali}

Pilih satu temuan dari Bagian~C.4, lalu tentukan dengan alasan: apakah baris
bermasalah itu sebaiknya diperbaiki, dibuang, atau dipertahankan apa adanya
dengan penanda. Sebutkan pihak mana di NusaMart yang berwenang memutuskan hal
itu. Jawaban tanpa alasan tidak memperoleh nilai penuh.

\section{Tugas Individu}

\subsection{Pekerjaan Wajib}
Lengkapi bus matrix awal untuk empat proses bisnis NusaMart dan jelaskan alasan
pemilihan grain setiap proses.

Bus matrix adalah tabel dengan proses bisnis sebagai baris dan dimensi sebagai
kolom; sebuah tanda diberikan bila proses tersebut memakai dimensi itu
\citep{kimball2013}. Empat proses yang harus dicakup: penjualan di kasir,
persediaan harian, cakupan promosi, dan retur barang. Format keluarannya
mengikuti Tabel~\ref{tab:m1-bus-matrix}.

\begin{table}[htbp]
\centering
\small
\caption{Kerangka bus matrix yang harus dilengkapi}
\label{tab:m1-bus-matrix}
\begin{tabular}{@{}p{0.30\textwidth}ccccc@{}}
\toprule
\textbf{Proses bisnis} & \textbf{Tanggal} & \textbf{Produk} & \textbf{Toko}
  & \textbf{Pelanggan} & \textbf{Promosi} \\
\midrule
Penjualan di kasir & & & & & \\
Persediaan harian  & & & & & \\
Cakupan promosi    & & & & & \\
Retur barang       & & & & & \\
\bottomrule
\end{tabular}
\end{table}

Kolom dimensi yang dipakai bersama oleh beberapa proses adalah calon
\istilah{conformed dimension} --- dimensi yang nantinya wajib memiliki definisi
tunggal di seluruh gudang data. Tandai mana saja yang termasuk, karena inilah
yang membuat perbandingan antarproses pada Modul~4 menjadi mungkin.

\subsection{Pertanyaan Analisis}

\begin{enumerate}[leftmargin=1.4em]
  \item Bila grain fakta penjualan ditetapkan pada tingkat struk, bukan pada
        tingkat baris struk, pertanyaan bisnis mana dari keempat yang Anda tulis
        yang menjadi tidak terjawab? Jelaskan mengapa.
  \item Kolom \perintah{harga\_jual} pada tabel \perintah{produk} berubah setiap
        kali harga disesuaikan, dan sistem operasional menimpa nilai lamanya.
        Apa akibatnya bagi laporan penjualan tahun lalu bila gudang data
        mengambil harga langsung dari tabel tersebut saat pelaporan?
  \item Transaksi berstatus batal masih tersimpan pada tabel
        \perintah{transaksi}. Menurut Anda, apakah baris tersebut layak dimuat
        ke gudang data? Uraikan konsekuensi kedua pilihan --- memuat dengan
        penanda, atau menyaringnya sejak awal.
\end{enumerate}

\section{Format Pengumpulan dan Checklist}

Kumpulkan inventaris sumber, desain konseptual versi pertama, dan bus matrix.
Pastikan jumlah baris berasal dari query aktual dan seluruh asumsi ditandai.

Seluruh berkas diletakkan pada direktori \berkas{modul-01-setup/} di repositori
pribadi, lalu di-\emph{push} paling lambat sebelum sesi laboratorium
berikutnya.

\begin{ceklis}
  \item \berkas{bukti-lingkungan.txt}
  \item \berkas{erd-sumber.png}
  \item \berkas{inventaris-sumber.md}
  \item \berkas{desain-konseptual-v1.md}
  \item \berkas{bus-matrix.md}
  \item \berkas{jawaban-analisis.md}
  \item Riwayat commit memperlihatkan pekerjaan bertahap, bukan satu commit
        berisi seluruh berkas.
\end{ceklis}

\section{Troubleshooting}
\label{sec:m1-troubleshooting}

\begin{table}[htbp]
\centering
\small
\caption{Masalah yang lazim muncul pada Modul 1 dan penanganannya}
\label{tab:m1-troubleshooting}
\begin{tabular}{@{}p{0.30\textwidth}p{0.62\textwidth}@{}}
\toprule
\textbf{Gejala} & \textbf{Penanganan} \\
\midrule
\emph{port is already allocated} &
  PostgreSQL lain memakai port tersebut. Ubah pemetaan port pada
  \berkas{docker-compose.yml} menjadi \perintah{"5443:5432"}, lalu sesuaikan
  koneksi DBeaver. Jangan mematikan PostgreSQL bawaan laptop. \\
Container berstatus \emph{exited} tepat setelah dinyalakan &
  Baca sebabnya dengan \perintah{docker compose logs postgres\_source}.
  Penyebab tersering adalah volume lama yang rusak; bersihkan dengan
  \perintah{docker compose down -v}, lalu \perintah{up} kembali. \\
\emph{password authentication failed} &
  Berkas \berkas{.env} belum dibuat dari \berkas{.env.example}, atau container
  masih memakai volume lama dengan kata sandi berbeda. \\
DBeaver \emph{connection refused} &
  Container sudah berjalan, tetapi port yang dituju keliru. Periksa kolom
  \emph{PORTS} pada \perintah{docker compose ps}. \\
Tabel ada tetapi kosong &
  Skrip seed belum selesai berjalan. Jalankan ulang
  \perintah{docker compose exec postgres\_source psql -U praktikum -d
  nusamart\_oltp -f /seed/nusamart\_kecil.sql}. \\
\perintah{relation "transaksi" does not exist} &
  Tabel berada pada schema \perintah{nusamart}, bukan \perintah{public}.
  Tulis nama lengkapnya, atau jalankan
  \perintah{SET search\_path TO nusamart;}. \\
\bottomrule
\end{tabular}
\end{table}

\section{Ringkasan}

Modul ini tidak menghasilkan satu tabel gudang data pun, dan memang tidak
seharusnya. Yang dihasilkan adalah dua hal yang menentukan mutu seluruh modul
berikutnya: pemahaman atas sumber yang didukung bukti, dan satu kalimat grain
yang dapat dipertanggungjawabkan.

Tiga gagasan perlu dibawa ke Modul~2. \emph{Pertama}, angka pada inventarisasi
harus berasal dari query; perkiraan tidak dapat dipakai untuk merekonsiliasi
hasil star query nanti. \emph{Kedua}, grain ditetapkan sebelum measure dan
dimensi dipilih --- membalik urutan ini menghasilkan tabel fakta yang tidak
dapat menjawab pertanyaan yang diinginkan. \emph{Ketiga}, temuan kualitas data
hari ini bukan sekadar catatan; ia adalah daftar pekerjaan yang akan
diselesaikan secara bertahap oleh baris \perintah{unknown} pada Modul~5,
standardisasi pada Modul~7, dan pengujian otomatis pada Modul~10.

Pada Modul~2, proses bisnis dan grain yang dirumuskan hari ini diterjemahkan
menjadi skema bintang pertama yang berjalan, lengkap dengan DDL, pemuatan data,
dan tiga star query yang hasilnya wajib sama persis dengan sumber.
